\chapter{Conclusion \& Limitations}
\label{sec:conclusion}

\section{Conclusion}

We study Open-Domain Large-Table Question Answering (OLTQA), where systems must retrieve relevant tables from a corpus and reason over tables whose full serialization may exceed the LLM context window.
To support public evaluation, we derive an OLTQA dataset from DataBench and evaluate table retrieval and end-to-end answer accuracy across four LLM backbones.
We further evaluate our method on a real-world financial corpus to examine practical constraints and domain-specific challenges.

Our method combines schema-level table representations for table retrieval with interactive multi-table reasoning.
On DataBench, it achieves the highest end-to-end answer accuracy among the evaluated OTQA methods across all four LLM backbones.
The retrieval results show that schema-level table representations are more effective than full-content table representations under the evaluated retrieval configurations.
However, high gold-table coverage alone does not guarantee strong end-to-end performance.
The interactive table access analysis highlights the importance of allowing the reasoner to obtain additional table information when the initial table context is insufficient.
Increasing the candidate size improves gold-table coverage, but the accuracy gains diminish while input-token usage continues to increase.

The real-world financial evaluation reveals additional challenges arising from similar table schemas, domain-specific semantics, and interaction with an operational database.
The closed-domain results show that providing the gold table does not eliminate reasoning errors, while the gap between closed-domain and open-domain performance highlights the additional difficulty of retrieving and identifying the appropriate table.
Overall, our results demonstrate the importance of addressing table retrieval and the reasoning stage jointly when developing OLTQA systems.

\section{Limitations}

Our study has several limitations.
First, our evaluation covers one public benchmark with $80$ tables and one financial corpus with $572$ tables.
Although these corpora contain large individual tables, they do not establish the effectiveness of our method over substantially larger table corpora.
The evaluated queries also have a single-table answer requirement.
Our multi-table reasoning design allows the reasoner to consider multiple retrieved candidate tables, but we do not evaluate queries that require combining information across tables or performing cross-table joins.

Second, schema-level table representations exclude cell values from table retrieval.
This design is effective on DataBench, but may miss relevant tables when the information needed to identify them is primarily contained in cell values.
This limitation is particularly relevant to financial tables with similar schemas.
Future work could incorporate selected cell-level information while controlling the cost of indexing and maintaining large, continuously updated tables.

Third, interactive multi-table reasoning introduces a trade-off between answer accuracy and inference cost.
Increasing the candidate size increases input-token usage without necessarily improving accuracy.
The reasoner may also generate invalid operations or repeatedly fail to recover from tool-execution errors.
Adaptive candidate selection, more efficient table context management, and improved error recovery are potential directions for improving efficiency and reliability.

Finally, the derived open-domain evaluation setting contains queries that may not uniquely identify their intended tables.
Automatic evaluation can also mark semantically valid answers as incorrect because of differences in answer formatting or numerical units.
These issues complicate the interpretation of answer accuracy.
Future evaluation could include clearer query specifications, annotations of alternative valid tables and answers, and more robust answer comparison.